%%%%%%%%%%  scan idx 116  |  folio 105  %%%%%%%%%%

\medskip
\noindent\textbf{Teorema 4.10.}
Sea $X_0$ un esquema afín de tipo finito sobre $k$.

\smallskip
\noindent\textbf{(a)}
Si $X_0$ es localmente una intersección completa, entonces
$\underline M_{X_0}$ es liso sobre $\underline C_\Lambda$.

\smallskip
\noindent\textbf{(b)}
Supongamos que, para un punto cerrado fijo $x\in X_0$, el abierto
$X_0-\{x\}$ es liso sobre $k$. Entonces ambos funtores
\[
\underline M_{X_0}
\longrightarrow
\underline M_{\operatorname{Spec}(\mathcal O_{X_0,x})}
\longrightarrow
\underline M_{\operatorname{Spec}(\widehat{\mathcal O}_{X_0,x})}
\]
son mínimamente lisos. En particular, si $(R,\mathfrak X)$ es una
deformación versal formal de $X_0$, entonces $(R,\mathfrak X_{(x)})$
y $(R,\mathfrak X_{(\widehat x)})$ son deformaciones versales
formales de $\operatorname{Spec}(\mathcal O_{X_0,x})$ y
$\operatorname{Spec}(\widehat{\mathcal O}_{X_0,x})$,
respectivamente.

\medskip
\noindent\textit{Demostración.}
Utilizamos la notación de la sección 3 para la cohomología de álgebras
conmutativas y escribimos $X_0=\operatorname{Spec}(A_0)$.

%%%%%%%%%%  scan idx 117  |  folio 106  %%%%%%%%%%

\smallskip
\noindent\textbf{(a)}
Basta mostrar que
$P:[\underline M_{X_0}]\longrightarrow\underline C_\Lambda$ es
liso. Sea $(R,\operatorname{Spec}(A))$ una deformación de
$\operatorname{Spec}(A_0)$ sobre
$R\in\operatorname{Ob}(\underline C_\Lambda)$. Dada una extensión
pequeña $R'\longrightarrow R$ en $\underline C_\Lambda$, de modo que
\[
0\longrightarrow k\longrightarrow R'\longrightarrow R
\longrightarrow0
\]
es exacta, debemos levantar $(R,\operatorname{Spec}(A))$ a $R'$. La
sucesión de transitividad para los $D^i$ con coeficientes en $A_0$,
aplicada a $R'\longrightarrow R\longrightarrow A$, da
\[
\cdots\longrightarrow D^1(A|R,A_0)
\longrightarrow D^1(A|R',A_0)
\longrightarrow D^1(R|R',k)\otimes_kA_0
\longrightarrow D^2(A|R,A_0)\longrightarrow\cdots .
\]
Como $R'\longrightarrow R$ es una extensión pequeña,
$D^1(R|R',k)$ es de dimensión uno sobre $k$, generado por la clase
$[R']$ representada por
$0\longrightarrow k\longrightarrow R'\longrightarrow R
\longrightarrow0$. Puesto que $A$ es plano sobre $R$,
\[
D^i(A|R,A_0)=D^i(A_0|k,A_0)
\qquad(i>0).
\]
En particular, $D^2(A|R,A_0)=0$ porque $A_0$ es localmente una
intersección completa, por la descripción en grados bajos de
3.12(8). Obtenemos, pues, una sucesión exacta
\[
\longrightarrow D^1(A_0|k,A_0)
\longrightarrow D^1(A|R',A_0)
\longrightarrow D^1(R|R',k)\otimes_kA_0
\longrightarrow0.
\]
Por tanto, existe una extensión $A'$ de $R'$-álgebras de $A$ por
$A_0$ cuya clase se aplica en $[R']\otimes1$. Equivalentemente,
$(R,\operatorname{Spec}(A))$ se levanta a la deformación
$(R',\operatorname{Spec}(A'))$.

\smallskip
\noindent\textbf{(b)}
Por 3.14 hay isomorfismos
\[
D^1(A_0|k,A_0)
\xrightarrow{\ \sim\ }
D^1((A_0)_x|k,(A_0)_x)
\xrightarrow{\ \sim\ }
D^1(\widehat{(A_0)_x}|k,\widehat{(A_0)_x}),
\]
y por tanto
\[
[\underline M_{X_0}(k[\varepsilon])]
\xrightarrow{\ \sim\ }
[\underline M_{\operatorname{Spec}(\mathcal O_{X_0,x})}
 (k[\varepsilon])]
\xrightarrow{\ \sim\ }
[\underline M_{\operatorname{Spec}(\widehat{\mathcal O}_{X_0,x})}
 (k[\varepsilon])].
\]
Por 2.7(a), resta mostrar que
\[
[\underline M_{X_0}]
\longrightarrow
[\underline M_{\operatorname{Spec}(\mathcal O_{X_0,x})}]
\longrightarrow
[\underline M_{\operatorname{Spec}(\widehat{\mathcal O}_{X_0,x})}]
\]
son funtores lisos.

Sea $(R,\operatorname{Spec}(A))$ una deformación de
$X_0=\operatorname{Spec}(A_0)$, y sea
$0\longrightarrow k\longrightarrow R'\longrightarrow R
\longrightarrow0$ una extensión pequeña. Como antes, los
homomorfismos $R'\longrightarrow R\longrightarrow A$ dan un diagrama
conmutativo exacto

%%%%%%%%%%  scan idx 118  |  folio 107  %%%%%%%%%%

\begin{center}
\scriptsize
\begin{tikzcd}[column sep=tiny,row sep=large]
{} \arrow[r]
& D^1(A_0|k,A_0) \arrow[r] \arrow[d,"\sim"']
& D^1(A|R',A_0) \arrow[r,"p"] \arrow[d]
& D^1(R|R',k)\otimes_kA_0 \arrow[r] \arrow[d]
& D^2(A_0|k,A_0) \arrow[r] \arrow[d,"\sim"]
& \cdots \\
{} \arrow[r]
& D^1(A_0|k,(A_0)_x) \arrow[r]
& D^1(A|R',(A_0)_x) \arrow[r]
& D^1(R|R',k)\otimes_k(A_0)_x \arrow[r]
& D^2(A_0|k,(A_0)_x) \arrow[r]
& \cdots \\
{} \arrow[r]
& D^1((A_0)_x|k,(A_0)_x) \arrow[r] \arrow[u]
& D^1(A_x|R',(A_0)_x) \arrow[r,"p_x"] \arrow[u]
& D^1(R|R',k)\otimes_k(A_0)_x \arrow[r] \arrow[u]
& D^2((A_0)_x|k,(A_0)_x) \arrow[r] \arrow[u]
& \cdots
\end{tikzcd}
\end{center}
Los isomorfismos indicados se deducen de que $A_0$ es liso sobre $k$
fuera del único punto cerrado $x$, de modo que
\[
D^i(A_0|k,(A_0)_x)
=D^i(A_0|k,A_0)_x
=D^i(A_0|k,A_0)
\qquad(i>0).
\]

Supongamos dada una flecha
\[
(R',\operatorname{Spec}(B))
\longrightarrow(R,\operatorname{Spec}(A_x))
\]
en $[\underline M_{\operatorname{Spec}((A_0)_x)}]$. Esto equivale a
una clase $[B]\in D^1(A_x|R',(A_0)_x)$ que satisface
$p_x([B])=[R']$. Una persecución en el diagrama da una clase
$[A']\in D^1(A|R',A_0)$ tal que $p([A'])=[R']$ y $A'_x=B$.
Así, $(R',\operatorname{Spec}(A'))$ es una deformación de
$\operatorname{Spec}(A_0)$ sobre $R'$ cuya localización en $x$ es
$\operatorname{Spec}(B)$. Esto demuestra que
\[
\underline M_{X_0}\longrightarrow
\underline M_{\operatorname{Spec}(\mathcal O_{X_0,x})}
\]
es liso.

La lisura de
\[
\underline M_{\operatorname{Spec}((A_0)_x)}
\longrightarrow
\underline M_{\operatorname{Spec}(\widehat{(A_0)_x})}
\]
se deduce del diagrama análogo y de las identidades
\[
D^i((A_0)_x|k,\widehat{(A_0)_x})
=\widehat{D^i((A_0)_x|k,(A_0)_x)}
=D^i((A_0)_x|k,(A_0)_x)
\qquad(i>0),
\]
puesto que $x$ es el único punto posiblemente no liso.
\hfill$\square$
